\section{Introduction}
\label{sec:intro}

Deployed LLM applications routinely place sensitive information in system prompts: API keys, internal URLs, proprietary configurations, and business logic.
Unlike traditional software, where secrets can be isolated in environment variables or vaults, LLMs have no mechanism to compartmentalize information within their context window.
Every token in the system prompt is, in principle, accessible to any sufficiently clever user query.
This creates a fundamental security gap: the very information that configures an LLM application is also the information most vulnerable to extraction~\cite{agarwal2024prompt,das2025system,wang2024raccoon}.

A growing body of work addresses this vulnerability through \emph{optimized} adversarial defenses.
Prompt Sensitivity Minimization~\cite{jawad2026psm} uses an LLM-as-optimizer loop to iteratively generate protective suffixes, reducing attack success rates to near-zero.
Defensive Prompt Patch~\cite{xiong2025dpp} evolves human-readable suffixes via hierarchical genetic algorithms.
Robust Prompt Optimization~\cite{zhou2024rpo} frames defense as minimax optimization over adversarial token sequences.
These approaches are effective but carry significant overhead: they require model-specific optimization, iterative refinement, and computational budgets that may be impractical for many deployments.

{\bf Can simple, handcrafted adversarial prompts provide meaningful protection?}
We test this question directly.
Rather than optimizing defensive text through expensive search procedures, we design three handcrafted adversarial ``firewalls'' of increasing strength---a security notice, a structured boundary marker, and a fake assistant completion---and place them immediately after sensitive fields within system prompts.
This placement exploits the well-documented recency bias in LLMs~\cite{liang2024prompts,chen2025dat}: models prioritize recent instructions, so a defensive prompt positioned after sensitive content can override extraction attempts that target that content.

We evaluate these firewalls on \gptmini across 65 diverse extraction attacks drawn from \raccoon~\cite{wang2024raccoon}, \citet{liang2024prompts}, and \citet{das2025system}, using 20 system prompts with three injected sensitive fields each.
Our best design---a structured boundary marker (\modfirewall)---reduces exact sensitive field leakage from 15.3\% to 6.4\%, a 58\% reduction ($p < 10^{-7}$) that matches the standard heuristic guardrail baseline.
The aggressive firewall (fake completion) also achieves significant protection (43\% reduction, $p < 10^{-3}$).

We also test an intriguing hypothesis: that firewalls might create \emph{selective permeability}, blocking general extraction while allowing access via exact-match queries that name specific sensitive values.
This hypothesis is partially refuted---firewalls suppress leakage broadly across query types, reducing exact-match leakage from 23\% to 0.5--4\% rather than preserving access for specific queries.

Our contributions are:
\begin{itemize}[leftmargin=*,itemsep=0pt,topsep=0pt]
    \item We propose and evaluate \emph{adversarial firewall prompts}: simple, handcrafted adversarial text placed after sensitive fields in system prompts, achieving up to 58\% leakage reduction without optimization or training.
    \item We conduct a controlled study across five defense conditions, six attack categories, and 5{,}650 total API calls, providing the first systematic comparison of handcrafted adversarial defenses for field-level protection within system prompts.
    \item We test and partially refute the selective permeability hypothesis, showing that adversarial firewalls suppress leakage broadly rather than creating query-type-dependent access control.
    \item We identify that injection-style attacks (\raccoon) account for nearly all successful extractions on \gptmini, while polite-request and diverse-strategy attacks produce near-zero leakage even without defenses.
\end{itemize}
